\documentclass{article}
\usepackage[osf]{mathpazo}
\usepackage{xspace}
\usepackage[margin=3cm]{geometry}
\usepackage{graphicx}
\graphicspath{{logo/}}
\setlength{\parindent}{0pt}
\setlength{\parskip}{1ex}
\pagenumbering{gobble}
\newcommand{\covid}{\textsc{covid-19}\xspace}
% ==============================================================================
\begin{document}
\begin{minipage}{0.5\linewidth}
  \includegraphics[height=0.8cm]{map-cuhs.eps}\\[1ex]
  \includegraphics[height=0.8cm]{smh.eps}     \\[1ex]
  \includegraphics[height=1.0cm]{uoft.eps}
\end{minipage}%
\begin{minipage}{0.5\linewidth}
  \begin{flushright}\small
    Sharmistha Mishra\\
    MAP Centre for Urban Health Solutions\\
    Li Ka Shing Knowledge Institute\\
    St Michael's Hospital,
    Unity Health Toronto\\
    University of Toronto\\
    sharmistha.mishra@utoronto.ca
  \end{flushright}
\end{minipage}
\\[3ex]
Dr.\ Yiming Shao \& Dr.\ Jianhong Wu\\
Editors-in-Chief of \textit{Infectious Disease Modelling}\\[2ex]
% ------------------------------------------------------------------------------
\textbf{Re. Submission of a research paper}\\[2ex]
% ------------------------------------------------------------------------------
Dear Editors,
\par
We are pleased to submit the attached manuscript entitled
\textit{Estimating effective reproduction number using
  generation time versus serial interval, with application to \covid
  in the Greater Toronto Area, Canada}
for consideration in \textit{Infectious Disease Modelling}.
\par
The effective reproduction number $R_e(t)$ is an essential metric
for monitoring the overall impact of interventions
in the ongoing \covid pandemic.
Popular methods of estimating $R_e(t)$ require only
an incidence time series and the generation time distribution.
However, since the generation time distribution is difficult to observe directly,
it is often approximated by the serial interval distribution.
While previous work has shown that such an approximation
can bias estimates of $R_e(t)$ due to exaggerated variance,
new challenges arise in the context of presymptomatic transmission, such as in \covid.
Namely, the generation time distribution is strictly positive,
while negative serial interval has been observed in
as much as 12\% of \covid transmissions.
\par
In this paper, we provide a method to recover the generation time distribution from
parametric definitions of the serial interval and incubation period distributions as inputs.
Unlike previous works, our method does not require person-level data, and could therefore use
pooled estimates of input parameter distributions obtained via meta-analyses.
We applied the method to infer the generation time distribution of \covid
using published serial interval and incubation period distributions.
We then compared estimated $R_e(t)$ for \covid in the Greater Toronto Area of Canada using:
negative-permitting versus non-negative serial interval distributions reported in the literature,
versus the inferred generation time distribution.
We found that the bias in $R_e(t)$ due to approximation of generation time with serial interval
was different for negative-permitting versus non-negative serial interval distributions.
Notably, whereas previous work showed that the approximation causes underestimation of $R_e(t)$,
we found the opposite was true for non-negative distributions fitted to serial interval data,
which have been used in several recent studies.
\par
To our knowledge, the proposed method to recover the generation time distribution
from \emph{parametric} definitions of the serial interval and incubation period distributions
is the first of its kind, and could be useful throughout and beyond the \covid pandemic.
Moreover, our finding that $R_e(t)$ can be \emph{overestimated} when approximating
the generation time distribution with the serial interval distribution in the context of \covid
runs counter to previous conclusions,
and thus has important implications regarding the severity of the epidemic
which are no-doubt relevant to the audience of \textit{Infectious Disease Modelling}.
\par
Thank you for your consideration, and we look forward to hearing from you.
\\[2ex]
Sincerely,\\[1ex]
Jesse Knight and Sharmistha Mishra
\end{document}
% ==============================================================================